\section{Setting and notation}
\label{sec:setting}

\subsection{The store}

We take the storage model of ClickHouse's MergeTree family as
given~\cite{schulze2024clickhouse,clickhousedocs}. The store the design is
intended for is Hanzo Datastore, a fork of ClickHouse~26.4 that keeps that
model~\cite{hanzodatastore}. \prior{} A table is a set of
parts $P_1,\dots,P_m$. A part is immutable: rows arrive in new parts and
background merges replace a set of parts with one part holding their union,
which is the log-structured merge discipline~\cite{oneil1996lsm}. Within a
part, rows are sorted by the ordering key and divided into granules, by default
of $8{,}192$ rows or of a byte budget, whichever binds first. The primary index
is sparse: it stores the ordering key of the first row of each granule, so key
predicates select ranges of granules rather than rows. Skipping indexes ---
minimum/maximum, set, Bloom filters over tokens or $n$-grams --- are written
inside the part, each summarising a run of granules, and exclude granules that
cannot satisfy a predicate; they are the modern form of small materialized
aggregates~\cite{moerkotte1998sma} and of the partitioning that data-skipping
work tunes deliberately~\cite{sun2014skipping}.

Two index types in that family concern retrieval directly. The text index
stores, per part, the rows in which each token occurs, with no weight, and is
merged rather than rebuilt when parts merge; it answers membership predicates.
The vector similarity index builds a USearch HNSW graph~\cite{malkov2020hnsw,
usearch} whose default granularity is large enough that in practice one graph
covers a whole part. Both are per part, so the work of a query is a sum over
the parts it visits.

\subsection{Rows, queries, scores}

\begin{definition}[Basis and representation]
A basis $\mathcal{B} = \bigcup_t \mathcal{B}_t$ is a set of elements, each
belonging to one type $t$: a word, a linked entity, a code symbol, a coordinate
of a learned sparse dictionary, a relation, or a structural motif. A row $x$
carries a sparse non-negative vector $c(x) \in \mathbb{R}_{\ge 0}^{\mathcal{B}}$
with few non-zero coordinates, and optionally a dense embedding
$v(x) \in \mathbb{R}^d$.
\end{definition}

\begin{definition}[Query and score]
A query is $Q = \sum_{i \in A(Q)} q_i B_i$ with non-negative $q_i$ and active
set $A(Q) \subseteq \mathcal{B}$, optionally with a dense vector $u(Q)$. The
score is
\[
  s(Q,x) \;=\; \sum_{i \in A(Q)} q_i\, c_i(x) \;+\; \lambda \langle u(Q), v(x)\rangle ,
  \qquad \lambda \ge 0 ,
\]
and the answer is the $k$ rows of highest score, written $S_k(Q)$, with
$\theta_k(Q)$ the $k$-th highest score.
\end{definition}

The form is the vector-space model~\cite{salton1975vsm} over a vocabulary that
is not only words, and the hybrid term is the usual linear combination of a
lexical and a dense leg. Restricting weights to be non-negative is what makes
the bounds in \S\ref{sec:routing} hold as stated; learned sparse encoders of
the SPLADE family produce non-negative weights by
construction~\cite{formal2021splade}.

\subsection{What cost means here}

\begin{definition}[Cost counters]
For a query $Q$ over a corpus of $N$ rows, $C_{\mathrm{parts}}$ is the number
of parts read, $C_{\mathrm{gran}}$ the number of granules read,
$C_{\mathrm{bytes}}$ the compressed bytes fetched, and $C_{\mathrm{route}}$ the
number of routing summaries examined. The store publishes the first three as
\texttt{SelectedParts}, \texttt{SelectedMarks} and
\texttt{ReadCompressedBytes}; the fourth has to be instrumented.
\end{definition}

\begin{definition}[Relevant neighbourhood]
$G^{*}(Q)$ is the set of granules containing at least one row of $S_k(Q)$, and
$P^{*}(Q)$ the set of parts containing at least one such granule.
\end{definition}

Any method that reads at granule granularity and returns $S_k(Q)$ exactly must
read every granule in $G^{*}(Q)$, so $|G^{*}(Q)|$ is a lower bound on
$C_{\mathrm{gran}}$ for exact retrieval, and the harness reports it as the
oracle row. It is a lower bound on reading the answer, not on finding it: a
method also pays for whatever it reads to establish that nothing else
qualifies, and that second term is what pruning and routing exist to reduce.

\paragraph{The hypothesis, stated so it can fail.}
Let $D$ be a query distribution and $\rho$ a recall target. Corpora of size $N$
are nested samples from one row distribution, so a row present at $N$ is
present at every larger $N$. The hypothesis is that, with the knobs at each $N$
set to the cheapest configuration reaching recall $\rho$,
\[
  \mathbb{E}_{Q \sim D}\!\left[C(Q;N)\right] \;=\; a N^{\alpha},
  \qquad \alpha < 1
\]
where $C$ is parts, granules or bytes touched per correct answer, and that
$\mathbb{E}[C_{\mathrm{gran}}] / \mathbb{E}[|G^{*}|]$ does not grow with $N$.
The second part matters more than the first: an exponent below one can also be
reached by a method whose answers get worse, which is why recall is held rather
than reported.

\begin{remark}[Where the hypothesis must fail]
Queries exist for which every method pays $\Theta(N)$: $k$ proportional to $N$;
a request that enumerates all matches of a basis element carried by every row;
a score distribution so flat that $\theta_k$ is met by most regions, leaving no
bound able to exclude one. The protocol runs such queries as a separate stratum
and reports them, rather than letting the average over a benign query
distribution stand for the design's behaviour.
\end{remark}
